\documentclass[10pt,a4paper]{amsart}
\usepackage[margin=19mm]{geometry}
\usepackage{amsmath,amssymb,mathtools}
\usepackage[scr=rsfs,cal=euler]{mathalfa}
\usepackage{xfrac}
\usepackage[unicode,hidelinks,bookmarks=false]{hyperref}
\newcommand{\ie}{i.e.\ }
\newcommand{\cf}{cf.\ }
\newcommand{\resp}{respectively\ }
\newcommand{\Subsection}{\subsection}
\providecommand{\nobreakdash}{\nobreak\hbox{-}}
\setlength{\emergencystretch}{2em}
\multlinegap0pt
\usepackage{fullpage}
\usepackage{latexsym,url,cases}
\numberwithin{equation}{section}
\usepackage{indentfirst}
\hypersetup{colorlinks=true,citecolor=blue,linkcolor=blue,urlcolor=blue}
\renewcommand{\baselinestretch}{1.4}
\renewcommand{\arraystretch}{1.2}
\newtheorem{theorem}{Theorem}[section] %
\newtheorem{lemma}[theorem]{Lemma} %
\newtheorem{corollary}[theorem]{Corollary} %
\newtheorem{proposition}[theorem]{Proposition} %
\newtheorem{problem}[theorem]{Problem}
\newtheorem*{problem1}{Problem 1}%
\newtheorem*{problem2}{Problem 2}%
\newtheorem{conjecture}[theorem]{Conjecture} %
\newtheorem{folklore}[theorem]{Folklore} %
\newtheorem{definition}[theorem]{Definition} %
\newtheorem{remark}[theorem]{Remark} %
\newtheorem{remark1}{Remark}%
\newtheorem{fact}{Fact} %
\newtheorem{claim}{Claim} %
\newtheorem*{que1}{Question 1}
\newtheorem*{que2}{Question 2}
\newtheorem*{theorem0}{Nathanson's theorem}
\newtheorem*{theoremA}{Theorem A}
\newtheorem*{theoremB}{Theorem B}
\newtheorem*{theoremC}{Theorem C}
\newtheorem*{theoremD}{Theorem D}
\newtheorem*{theoremE}{Theorem E}
\newtheorem*{theoremF}{Theorem F}
\newtheorem*{theoremG}{Theorem G}
\newtheorem*{theoremH}{Theorem H}
\begin{document}
\title[On a problem of minimal additive complements for not eventua]{On a problem of minimal additive complements for not eventually periodic $S$-difference sets}
\author{Min Tang; Wenjing He}
\date{}
\maketitle
\noindent\emph{Source and licence.} On a problem of minimal additive complements for not eventually periodic $S$-difference sets. Version arXiv:2607.27059v1 (CC BY 4.0). This material is distributed under the Creative Commons Attribution 4.0 International licence, http://creativecommons.org/licenses/by/4.0/, as recorded in the arXiv OAI licence metadata for this exact version and shown on the abstract page https://arxiv.org/abs/2607.27059v1. Full rights notice: licenses/arxiv-2607.27059-CC-BY-4.0.txt.\par\smallskip
\noindent\textit{Editorial scope.} Counted unit: the original \texttt{lemma} environment \texttt{lem2} at source lines 153--155 together with its complete proof at lines 157--180; the delivered document keeps source lines 144--181 so that every referenced label and every citation resolves inside it. Native editable source is the paper's own concatenated TeX; the AMS wrapper is editorial formatting only. No publisher class, upstream script or unrelated artwork is redistributed.
 \section{Lemmas}

 \begin{lemma}\label{lem1} (See \cite[Theorem 1]{Brauer42}.) Let $a_1 \leqslant a_2 \leqslant \dots \leqslant a_k$ be relatively prime positive integers. We put
$$S = S(a_1, a_2, \dots, a_k) = a_2 + a_3 + \dots + a_{k-1} + a_1 a_k.$$
For $n > S$ the Diophantine equation
$$a_1 x_1 + a_2 x_2 + \dots + a_k x_k = n$$
always has solutions in positive integers $x_1, x_2, \dots, x_k$.
\end{lemma}


 \begin{lemma}\label{lem2} Let $d,k\geqslant 2$ be positive integers and $S=\{ds_1,\ldots,ds_k\}$ be a set of positive integers with $\gcd(s_1,\ldots,s_k)=1$.
 If $W = \{ w_1 < w_2 < \cdots \} \subseteq \mathbb{N}$ is an $\{s_1,\ldots,s_k\}$-difference set and $W$ has a minimal additive complement, then $dW$ is an $S$-difference set and $dW$ has a minimal additive complement.
  \end{lemma}

\begin{proof} It is easy to see that if $W$ is an $\{s_1,\ldots,s_k\}$-difference set, then $d W$ is an $S$-difference set.

Suppose that $C$ is a minimal additive complement to $W$.
 Let
$$C_d=\bigcup_{r=0}^{d-1}(dC + r).$$
We shall show that $C_d$ is a minimal additive complement to $d W$.

First, we prove that $C_d$ is an additive complement to $d W$.
For any $n \in \mathbb{Z}$, write $n= dq+r$, where $0 \leq r < d$. Since
$C+W=\mathbb{Z}$, there exist $w \in W$, $c \in C$ such that $w+c=q$.
Thus
$$n=dq+r=dw+dc+r\in d W+C_d.$$

Next, we show that for any $x \in C_d$, $d W+(C_d\setminus\{x\})\neq \mathbb{Z}.$

Write $x=dc_0+r_0$ for some $c_0 \in C$ and $0 \leq r_0 < d$. Since
$W+(C\setminus\{c_0\})\neq \mathbb{Z}$, there exists an integer $q_0\notin W+(C\setminus\{c_0\})$.
Then $$dq_0+r_0\notin d W+(C_d\setminus\{x\}).$$
Otherwise, if there exist $w'\in W$, $c'\in C$ and $0\leqslant r'\leqslant d-1$ such that $$dq_0+r_0=dw'+dc'+r',$$
then $w'+c'=q_0$ and $r'=r_0.$
It follows that $c'\neq c_0$, which contradicts that $q_0\notin W+(C\setminus\{c_0\})$.

  This completes the proof of Lemma \ref{lem2}.
\end{proof}


\begin{thebibliography}{99}
\bibitem[Brauer42]{Brauer42} A. Brauer, {\it On a problem of partitions}, Amer. J. Math. 64(1942), 299-312.
\end{thebibliography}
\end{document}
